% Options for packages loaded elsewhere
\PassOptionsToPackage{unicode}{hyperref}
\PassOptionsToPackage{hyphens}{url}
\documentclass[
  ignorenonframetext,
]{beamer}
\newif\ifbibliography
\usepackage{pgfpages}
\setbeamertemplate{caption}[numbered]
\setbeamertemplate{caption label separator}{: }
\setbeamercolor{caption name}{fg=normal text.fg}
\beamertemplatenavigationsymbolsempty
% remove section numbering
\setbeamertemplate{part page}{
  \centering
  \begin{beamercolorbox}[sep=16pt,center]{part title}
    \usebeamerfont{part title}\insertpart\par
  \end{beamercolorbox}
}
\setbeamertemplate{section page}{
  \centering
  \begin{beamercolorbox}[sep=12pt,center]{section title}
    \usebeamerfont{section title}\insertsection\par
  \end{beamercolorbox}
}
\setbeamertemplate{subsection page}{
  \centering
  \begin{beamercolorbox}[sep=8pt,center]{subsection title}
    \usebeamerfont{subsection title}\insertsubsection\par
  \end{beamercolorbox}
}
% Prevent slide breaks in the middle of a paragraph
\widowpenalties 1 10000
\raggedbottom
\AtBeginPart{
  \frame{\partpage}
}
\AtBeginSection{
  \ifbibliography
  \else
    \frame{\sectionpage}
  \fi
}
\AtBeginSubsection{
  \frame{\subsectionpage}
}
\usepackage{iftex}
\ifPDFTeX
  \usepackage[T1]{fontenc}
  \usepackage[utf8]{inputenc}
  \usepackage{textcomp} % provide euro and other symbols
\else % if luatex or xetex
  \usepackage{unicode-math} % this also loads fontspec
  \defaultfontfeatures{Scale=MatchLowercase}
  \defaultfontfeatures[\rmfamily]{Ligatures=TeX,Scale=1}
\fi
\usepackage{lmodern}
\ifPDFTeX\else
  % xetex/luatex font selection
\fi
% Use upquote if available, for straight quotes in verbatim environments
\IfFileExists{upquote.sty}{\usepackage{upquote}}{}
\IfFileExists{microtype.sty}{% use microtype if available
  \usepackage[]{microtype}
  \UseMicrotypeSet[protrusion]{basicmath} % disable protrusion for tt fonts
}{}
\makeatletter
\@ifundefined{KOMAClassName}{% if non-KOMA class
  \IfFileExists{parskip.sty}{%
    \usepackage{parskip}
  }{% else
    \setlength{\parindent}{0pt}
    \setlength{\parskip}{6pt plus 2pt minus 1pt}}
}{% if KOMA class
  \KOMAoptions{parskip=half}}
\makeatother
\usepackage{color}
\usepackage{fancyvrb}
\newcommand{\VerbBar}{|}
\newcommand{\VERB}{\Verb[commandchars=\\\{\}]}
\DefineVerbatimEnvironment{Highlighting}{Verbatim}{commandchars=\\\{\}}
% Add ',fontsize=\small' for more characters per line
\newenvironment{Shaded}{}{}
\newcommand{\AlertTok}[1]{\textcolor[rgb]{1.00,0.00,0.00}{\textbf{#1}}}
\newcommand{\AnnotationTok}[1]{\textcolor[rgb]{0.38,0.63,0.69}{\textbf{\textit{#1}}}}
\newcommand{\AttributeTok}[1]{\textcolor[rgb]{0.49,0.56,0.16}{#1}}
\newcommand{\BaseNTok}[1]{\textcolor[rgb]{0.25,0.63,0.44}{#1}}
\newcommand{\BuiltInTok}[1]{\textcolor[rgb]{0.00,0.50,0.00}{#1}}
\newcommand{\CharTok}[1]{\textcolor[rgb]{0.25,0.44,0.63}{#1}}
\newcommand{\CommentTok}[1]{\textcolor[rgb]{0.38,0.63,0.69}{\textit{#1}}}
\newcommand{\CommentVarTok}[1]{\textcolor[rgb]{0.38,0.63,0.69}{\textbf{\textit{#1}}}}
\newcommand{\ConstantTok}[1]{\textcolor[rgb]{0.53,0.00,0.00}{#1}}
\newcommand{\ControlFlowTok}[1]{\textcolor[rgb]{0.00,0.44,0.13}{\textbf{#1}}}
\newcommand{\DataTypeTok}[1]{\textcolor[rgb]{0.56,0.13,0.00}{#1}}
\newcommand{\DecValTok}[1]{\textcolor[rgb]{0.25,0.63,0.44}{#1}}
\newcommand{\DocumentationTok}[1]{\textcolor[rgb]{0.73,0.13,0.13}{\textit{#1}}}
\newcommand{\ErrorTok}[1]{\textcolor[rgb]{1.00,0.00,0.00}{\textbf{#1}}}
\newcommand{\ExtensionTok}[1]{#1}
\newcommand{\FloatTok}[1]{\textcolor[rgb]{0.25,0.63,0.44}{#1}}
\newcommand{\FunctionTok}[1]{\textcolor[rgb]{0.02,0.16,0.49}{#1}}
\newcommand{\ImportTok}[1]{\textcolor[rgb]{0.00,0.50,0.00}{\textbf{#1}}}
\newcommand{\InformationTok}[1]{\textcolor[rgb]{0.38,0.63,0.69}{\textbf{\textit{#1}}}}
\newcommand{\KeywordTok}[1]{\textcolor[rgb]{0.00,0.44,0.13}{\textbf{#1}}}
\newcommand{\NormalTok}[1]{#1}
\newcommand{\OperatorTok}[1]{\textcolor[rgb]{0.40,0.40,0.40}{#1}}
\newcommand{\OtherTok}[1]{\textcolor[rgb]{0.00,0.44,0.13}{#1}}
\newcommand{\PreprocessorTok}[1]{\textcolor[rgb]{0.74,0.48,0.00}{#1}}
\newcommand{\RegionMarkerTok}[1]{#1}
\newcommand{\SpecialCharTok}[1]{\textcolor[rgb]{0.25,0.44,0.63}{#1}}
\newcommand{\SpecialStringTok}[1]{\textcolor[rgb]{0.73,0.40,0.53}{#1}}
\newcommand{\StringTok}[1]{\textcolor[rgb]{0.25,0.44,0.63}{#1}}
\newcommand{\VariableTok}[1]{\textcolor[rgb]{0.10,0.09,0.49}{#1}}
\newcommand{\VerbatimStringTok}[1]{\textcolor[rgb]{0.25,0.44,0.63}{#1}}
\newcommand{\WarningTok}[1]{\textcolor[rgb]{0.38,0.63,0.69}{\textbf{\textit{#1}}}}
\setlength{\emergencystretch}{3em} % prevent overfull lines
\providecommand{\tightlist}{%
  \setlength{\itemsep}{0pt}\setlength{\parskip}{0pt}}
% Auto-generated by infrastructure/rendering/_pdf_latex_helpers.py::extract_math_font_preamble.
% Minimal header passed to Pandoc via -H so Beamer slides pick up the same math font
% as the combined-PDF path. Adjust the manuscript's preamble.md to change the font globally.
\usepackage{unicode-math}
\setmathfont{latinmodern-math.otf}

% Slide layout defaults for warning-clean scientific decks.
\usepackage{etoolbox}
\IfFileExists{xurl.sty}{\usepackage{xurl}}{}
\IfFileExists{seqsplit.sty}{\usepackage{seqsplit}}{\newcommand{\seqsplit}[1]{#1}}
\protected\def\breakseq#1{\seqsplit{#1}}
\protected\def\breaktt#1{\begingroup\ttfamily\seqsplit{#1}\endgroup}
\setlength{\emergencystretch}{6em}
\tolerance=5000
\hbadness=10000
\hfuzz=1pt
\setlength{\tabcolsep}{2pt}
\AtBeginEnvironment{longtable}{\tiny\renewcommand{\arraystretch}{0.86}\setlength{\tabcolsep}{1pt}}
\AtBeginEnvironment{tabular}{\tiny\renewcommand{\arraystretch}{0.86}\setlength{\tabcolsep}{1pt}}

% Natbib and cross-reference fallbacks — slides don't load natbib
% or cleveref, but combined-PDF manuscript prose may emit these
% commands. The fallback renders citations as a bracketed key list
% and cross-references as detokenized labels so slides don't fail on
% undefined control sequences or raw underscores. \providecommand is
% a no-op if packages load later.
\providecommand{\citep}[1]{[#1]}
\providecommand{\citet}[1]{#1}
\providecommand{\citealp}[1]{#1}
\providecommand{\citeauthor}[1]{#1}
\providecommand{\citeyear}[1]{#1}
\providecommand{\cref}[1]{\texttt{\detokenize{#1}}}
\providecommand{\Cref}[1]{\texttt{\detokenize{#1}}}

\newtheorem{proposition}{Proposition}
\newtheorem{hypothesis}{Hypothesis}
\usepackage{bookmark}
\IfFileExists{xurl.sty}{\usepackage{xurl}}{} % add URL line breaks if available
\urlstyle{same}
% fallback for those not using the hyperref driver hyperxmp:
\makeatletter
\@ifundefined{xmpquote}{\newcommand{\xmpquote}[1]{#1}}{}
\makeatother
\hypersetup{
  hidelinks,
  pdfcreator={LaTeX via pandoc}}

\author{\texorpdfstring{}{}}
\date{}

\begin{document}

\section{Reproducibility and Provenance}\label{sec:reproducibility}

\begin{frame}[fragile,allowframebreaks]{Reproducibility and Provenance}
The reproducibility contract is:

\begin{enumerate}
\tightlist
\item
  construct the same frozen parameter dataclass;
\item
  use the same generator ID and the same deterministic seed;
\item
  compare the canonical artifact digest;
\item
  record the parameter schema, typed parameters, taxonomy, evidence
  boundary, canonical byte serialization, objective metrics, and
  encoding request;
\item
  if a file is written, inspect the decoded media and compare
  dimensions, channels, rates, frame counts, durations, timing offsets,
  and encoded hash;
\item
  retain the v2 manifest and source-data sidecars with the generated
  figure or table.
\end{enumerate}

The contract is deliberately stronger than ``the file can be
downloaded.'' Research-software guidance treats versioned code,
executable procedures, dependencies, and retained inputs as part of the
reproducible object \breakseq{[@sandve2013simple; @wilson2014bestpractices]}.
FAIR guidance further asks that digital research objects be findable,
accessible, interoperable, and reusable, while noting that software has
lifecycle and maintenance constraints that are not identical to those of
static data \breakseq{[@wilkinson2016fair; @lamprecht2020fairsoftware]}.
DuckRabbit operationalizes those principles with typed metadata,
deterministic hashes, source-data sidecars, explicit licenses, and
release gates rather than by implying that a generated media file is a
complete scientific replication.

The package reports implemented, planned, input\_required as its catalog
status vocabulary and currently covers audio\_visual, auditory, visual.
The evidence layer contains 18 entry records backed by 36 source
records. The publication workflow produces 15 figures and 10 tables from
the live registry.

The public software record is
\href{https://github.com/docxology/DuckRabbit}{the DuckRabbit public
GitHub repository}, with Daniel Ari Friedman of the Active Inference
Institute as author. The release is archived on Zenodo at
\href{https://doi.org/10.5281/zenodo.21419693}{10.5281/zenodo.21419693}
(concept DOI, always resolves to the latest version), with
\href{https://doi.org/10.5281/zenodo.21419694}{10.5281/zenodo.21419694}
identifying this v0.5.0 release specifically.

The canonical buffer is the primary identity of a stimulus. Encoded
files are delivery artifacts with format-specific tolerances. The
observer protocol is a separate evidence layer: a manifest can prove
what was presented without proving what an observer experienced.

Synthetic psychophysics is a third, deliberately bounded object. The
model identity, version, feature schema, weights, temperature, seed,
calibration statement, canonical reference/comparison digests, and
\breaktt{human\_data=false} flag are retained in
\breaktt{output/data/synthetic\_psychophysics.json}. Re-running this
diagnostic checks deterministic model orchestration and stimulus-feature
sensitivity; it does not estimate a human observer or replace empirical
psychophysics.

The v2 verifier is a relational check rather than a field-presence
check. It rejects summaries whose dimensions, dtype, byte count, rates,
duration, signal range, frame deltas, or signed audiovisual offset
disagree. Public manifest and inspection mappings are recursively frozen
after validation, and rational frame rates are reduced before entering
the clock contract. The read-only capability probe records whether
Pillow, NumPy, ffmpeg, or ffprobe are available in the current
environment; unavailable delivery paths remain explicit rather than
being silently downgraded.

The minimal regeneration commands are:

\begin{Shaded}
\begin{Highlighting}[]
\ExtensionTok{uv}\NormalTok{ run pytest }\AttributeTok{{-}q} \AttributeTok{{-}{-}cov}\OperatorTok{=}\NormalTok{src/duckrabbit }\AttributeTok{{-}{-}cov{-}fail{-}under}\OperatorTok{=}\NormalTok{90}
\ExtensionTok{uv}\NormalTok{ run python scripts/generate\_publication\_outputs.py}
\ExtensionTok{uv}\NormalTok{ run python scripts/z\_generate\_manuscript\_variables.py}
\end{Highlighting}
\end{Shaded}
\end{frame}

\begin{frame}[fragile,allowframebreaks]{Data availability and software
citation}
\protect\phantomsection\label{data-availability-and-software-citation}
No participant records, fitted observer coefficients, or human
psychophysics are bundled. The synthetic diagnostic is explicitly model
output with \breaktt{human\_data=false} and \breaktt{training\_data=none}.
Reuse should cite Daniel Ari Friedman and the release metadata in
\texttt{CITATION.cff}, together with the source-specific scholarship in
\texttt{references.bib}. Software-citation principles recommend citing
the software object itself, with enough version and identity information
to distinguish one release from another \breakseq{[@smith2016software]}. The
DOI field carries the real, minted concept DOI
(\breaktt{10.5281/zenodo.21419693}) and \texttt{doi\_status} is
\texttt{published}; no resolver URL was manufactured before that
identifier existed.
\end{frame}

\end{document}
